% FILE: 02_lineage_and_context.tex
% Project: research-prompt-manufactory
% Thesis Role: prompt-manufacturing-and-subagent-command-surface

\section{Lineage and Context}
\label{sec:research-prompt-manufactory-lineage}

\subsection{2016 Language-Model Lesson}
\label{sec:research-prompt-manufactory-lineage-2016}

The thesis lineage begins with a 2016 experiment in local language-model text imitation:
a model trained to reproduce text does not conserve the consequences of that text. A
language model that has seen a conformance-checking specification does not thereby
implement conformance checking. The gap between linguistic surface and operational
consequence is irreducible by training alone.

The Prompt Manufactory is a direct descendant of this lesson applied to the command
surface of multi-agent systems. A hand-written prompt that describes a subagent's role
in natural language is, at its core, a text-imitation act: it states a behavior without
binding that behavior to any formal structure that could be audited, replayed, or
compared against a declared law. The 2016 lesson predicts that such prompts will drift
from their intended semantics as the research program evolves, because the prompts have
no formal tether to the ontology they claim to reflect.

The Prompt Manufactory resolves this by treating prompt manufacture as a formal
production act. The subagent warrant for a Census role is not written to describe what
Census agents do; it is manufactured from the \texttt{pm:SubagentRole} instance in
\texttt{subagent-role-law.ttl} that defines what Census agents are permitted to inspect,
what surfaces they are forbidden from touching, and what output contract they must
satisfy. The warrant is a rendering of law, not a description of intent.

\subsection{Enterprise Process Gap}
\label{sec:research-prompt-manufactory-lineage-enterprise}

The enterprise process gap identified in the thesis lineage is the absence of a
verifiable bridge between declared process models and actual executed behavior. In
enterprise settings, process documentation is authored separately from execution
infrastructure. Process maps describe what should happen; execution logs record what
did happen. The gap between them is the terrain where conformance defects accumulate
invisibly.

The same gap exists in multi-agent AI systems. A research workflow documented in a
CLAUDE.md or a hand-written prompt file is process documentation with no execution
binding. When a subagent executes, it reads the prompt and produces behavior; but the
prompt itself carries no receipt, no provenance, and no derivation from the formal
process model. Whether the subagent's actual behavior conforms to the declared workflow
is unverifiable from the prompt alone.

The Prompt Manufactory closes this gap at the command-surface layer. By manufacturing
warrants from the same ontology that governs the workflow (specifically,
\texttt{workflow-law.ttl} defining INTEL and REMEDIATE workflow structures with
8 phases each), the Manufactory creates a traceable chain from declared law to
operational instruction. A subagent executing under a manufactured warrant is executing
under a command that is provably derived from the workflow law it is supposed to
implement.

\subsection{Relationship to the Chatman Equation}
\label{sec:research-prompt-manufactory-lineage-equation}

% AUTHOR THESIS / INTERPRETATION
The Chatman Equation $A = \mu(O^*)$ asserts that an artifact is the manufacturing
function $\mu$ applied to an ontology $O^*$. In this framing, the Prompt Manufactory
instantiates the equation at the command-surface layer: each manufactured warrant $A_i$
is $\mu_i(O^*)$, where $\mu_i$ is the composition of a SPARQL SELECT query with a Tera
template, and $O^*$ is the union of the 8 ontology files loaded into the RDF graph.

The critical property the equation confers is that the manufacturing function $\mu$ is
inspectable. Given the ontology hash recorded in
\texttt{ggen/.ggen/sync-state.json} (\texttt{01b52b81...}), the manifest hash, and the
generation rules in \texttt{ggen.toml}, any auditor can reconstruct the manufacturing
pipeline and verify that a given warrant $A_i$ is indeed $\mu_i(O^*)$ and not a
hand-written artifact inserted into the emission directory without provenance.

This reproducibility property is what distinguishes the Manufactory from a prompt
library. A prompt library stores artifacts; the Manufactory stores the manufacturing
function and emits artifacts as its outputs. The artifacts are replaceable; the function
is the authority.

\subsection{Relationship to Receipts and Replay}
\label{sec:research-prompt-manufactory-lineage-receipts}

The receipt chain is the evidentiary spine of the Process Intelligence research program.
Every manufacturing act that produces a consequential artifact must emit a receipt ---
a proof-of-execution record that can later be replayed to verify the artifact's
provenance. The Prompt Manufactory implements this doctrine at the warrant layer.

The \texttt{prompt-receipt-ledger.md} documents the receipt chain for each emitted
warrant. For the proof-specimen warrant
\texttt{PI\_RESEARCH\_PROGRAM\_INTEL\_001.md}, the ledger records: source ontology
(\texttt{research-program-law.ttl}), source instance URI
(\texttt{https://pi-research.dev/programs\#PI\_RESEARCH\_PROGRAM\_INTEL\_001}),
manufacturing timestamp (\texttt{2026-06-01T00:00:00Z}), manufacture method (fallback
manual due to ggen v5 template rendering failure on empty SPARQL result set), and
proof status (PARTIAL).

The \texttt{.ggen/receipts/phase-2-index.json} extends this to the batch level:
42 artifacts manufactured across 5 artifact types, with a ledger reference to
the full RECEIPT\_LEDGER for the manufacturing phase. The manufacturing manifest
\texttt{manufacturing-manifest-20260601.yaml} independently records 41 artifacts
emitted across 6 rules with receipts-equal-artifacts gate PASS --- a 1-artifact
discrepancy with the phase-2-index that has not been reconciled (see open questions).

Together these receipts form a replay surface: given the receipts, the manufacturing
pipeline can be re-executed to verify that the emitted artifacts match their declared
provenance. This is the receipt-and-replay doctrine instantiated at the prompt layer.
